\documentclass[11pt]{article}

\usepackage[margin=1in]{geometry}
\usepackage{amsmath, amssymb, amsthm}
\usepackage{booktabs}
\usepackage{mathtools}
\usepackage{hyperref}

\newtheorem{remark}{Remark}
\newcommand{\R}{\mathbb{R}}
\newcommand{\Om}{\Omega}
\newcommand{\eps}{\varepsilon}
\newcommand{\ustar}{u^\star}
\newcommand{\E}{\mathbb{E}}
\DeclareMathOperator*{\argmin}{arg\,min}

\title{The infinite-data limit in physics-informed operator learning:\\
which error vanishes, and why it is free}
\author{Marius Zeinhofer}
\date{\today}

\begin{document}
\maketitle

\begin{abstract}
Physics-informed operator learning trains on a \emph{residual} loss that requires no
solution labels. Consequently a fresh training sample costs only the draw of a random
right-hand side---no PDE solve---and the \emph{infinite-data limit}, in which no sample is
ever seen twice, is a practical operating point rather than a thought experiment. We
recall the classical decomposition of the excess risk into approximation, estimation
(generalization), and optimization errors, and observe that one-pass training on freshly
drawn samples eliminates the estimation error \emph{exactly}: stochastic gradient descent
then acts directly on the population risk. The remaining errors are the approximation
error of the model class and the optimization error of the finite compute budget. This
note fixes the terminology, states the prediction for a ``test error versus dataset
size'' figure at fixed optimization budget, and records the experimental-design caveats.
\end{abstract}

\section{Setting: two risks, no labels}

Let $\mu$ be a probability distribution over right-hand sides $f$ (in our experiments:
random coefficients on $K$ Fourier modes), and let $u_\theta$ denote the neural operator
(an FNO) with parameters $\theta\in\Theta$, mapping $f$ to an approximate solution of
\begin{equation}
  -\Delta u = f \quad\text{in }\Om, \qquad u=0 \quad\text{on }\partial\Om .
\end{equation}
After discretization ($Q_1$ finite elements, stiffness matrix $A$, load vector $b(f)$;
see the companion note on the transition preconditioner), the \emph{physics-informed}
per-sample loss is a residual functional, e.g.\ the preconditioned least-squares loss
\begin{equation}
  \ell(\theta; f) \;=\; \tfrac12\,\bigl\| P\bigl(A\,u_\theta(f) - b(f)\bigr) \bigr\|^2 ,
  \label{eq:pls}
\end{equation}
or the (preconditioned) Deep Ritz energy. The object of interest is the
\emph{population risk} and its empirical counterpart on a dataset
$D_n=\{f_1,\dots,f_n\}\overset{\text{iid}}{\sim}\mu$:
\begin{equation}
  R(\theta) \;=\; \E_{f\sim\mu}\,\ell(\theta;f),
  \qquad
  R_n(\theta) \;=\; \frac1n \sum_{i=1}^n \ell(\theta;f_i).
\end{equation}
The evaluation metric (relative $L^2$ error against the manufactured analytic solution;
see the remark below) is a different functional, but the discussion is insensitive to
this distinction: what matters is that training only ever touches $\ell$, which does not
contain the solution $\ustar(f)$.

\begin{remark}[Why streaming is free here]\label{rem:free}
Evaluating \eqref{eq:pls} requires $b(f)$---an assembly---but \emph{not} $\ustar(f)$. In
our benchmark this asymmetry is softened: the dataset uses \emph{manufactured solutions}
(the random coefficients define an analytic $\ustar$, and $f=-\Delta\ustar$ in closed
form), so labels also cost nothing and a fresh sample is a random-number draw plus the
evaluation of two closed-form fields---for \emph{every} loss. The asymmetry the
experiment ultimately speaks to is the general case, where $f\sim\mu$ is prescribed and
$\ustar(f)$ requires a numerical solve: there, \emph{data-driven} operator learning
prices the infinite-data limit at one solver call per gradient contribution, while the
physics-informed loss keeps it free. Within the benchmark, streaming simply lets us
isolate the estimation error cleanly, at no extra cost, for whichever loss we train.
\end{remark}

\section{The error decomposition}

Following the classical large-scale-learning decomposition (Bottou \& Bousquet, 2008),
write $R^\star=\inf_u R(u)$ for the risk of the best measurable predictor,
$\theta_{\mathrm{app}}\in\argmin_{\theta\in\Theta} R(\theta)$ for the best operator in
the model class, $\theta_n\in\argmin_{\theta\in\Theta} R_n(\theta)$ for the empirical
risk minimizer, and $\tilde\theta$ for what the optimizer actually returns after a
compute budget of $T$ gradient evaluations. Then the excess risk splits as
\begin{equation}
  \underbrace{R(\tilde\theta)-R^\star}_{\text{excess risk}}
  \;=\;
  \underbrace{R(\theta_{\mathrm{app}})-R^\star}_{\eps_{\mathrm{app}}\;\text{(approximation)}}
  \;+\;
  \underbrace{R(\theta_n)-R(\theta_{\mathrm{app}})}_{\eps_{\mathrm{est}}\;\text{(estimation)}}
  \;+\;
  \underbrace{R(\tilde\theta)-R(\theta_n)}_{\eps_{\mathrm{opt}}\;\text{(optimization)}} .
  \label{eq:decomp}
\end{equation}
The three terms are controlled by three different resources: $\eps_{\mathrm{app}}$ by
model capacity (FNO width, modes, depth), $\eps_{\mathrm{est}}$ by the dataset size $n$
(typically $\eps_{\mathrm{est}}=O((C_\Theta/n)^\alpha)$ with a complexity measure
$C_\Theta$ and $\alpha\in[\tfrac12,1]$), and $\eps_{\mathrm{opt}}$ by the budget $T$ and
by the \emph{conditioning of the loss}---which is where the preconditioner $P$ of the
companion note enters.

\section{Streaming eliminates the estimation error exactly}

Suppose each optimization step draws a \emph{fresh} minibatch
$f_{t,1},\dots,f_{t,B}\overset{\text{iid}}{\sim}\mu$, never reusing a sample
(``streaming'' or \emph{one-pass} training). Then for every $t$,
\begin{equation}
  \E\Bigl[\frac1B\sum_{j=1}^B \nabla_\theta\,\ell(\theta_t; f_{t,j})\Bigr]
  \;=\; \nabla R(\theta_t),
\end{equation}
i.e.\ the update is an unbiased stochastic gradient of the \emph{population} risk. No
finite empirical risk $R_n$ ever appears; the algorithm is stochastic approximation
(Robbins--Monro) applied directly to $R$. The decomposition \eqref{eq:decomp} collapses to
\begin{equation}
  R(\tilde\theta)-R^\star \;=\; \eps_{\mathrm{app}} + \eps_{\mathrm{opt}},
  \qquad \eps_{\mathrm{est}} \equiv 0 .
\end{equation}
So yes: the streaming limit removes one of the three error terms \emph{identically}, not
merely asymptotically. Two consequences and one honest caveat:

\begin{enumerate}
  \item \textbf{No generalization gap.} The running training loss is itself an unbiased
  estimate of the test loss; overfitting is impossible in the usual sense, and early
  stopping is unnecessary (it would only truncate the budget).
  \item \textbf{No memorization.} With a finite dataset and $T\gg n/B$ the optimizer
  spends most of its budget re-fitting samples it has already fit, driving
  $R_n(\tilde\theta)$ below $R(\tilde\theta)$; the streaming run spends every gradient
  evaluation on new information.
  \item \textbf{Caveat: the sampling noise does not disappear---it moves.} The variance
  of the minibatch gradient is the same whether the batch is fresh or recycled, and it
  enters the convergence rate of stochastic approximation. Eliminating
  $\eps_{\mathrm{est}}$ does not mean finite-sample effects vanish; it means they are
  accounted for \emph{inside} $\eps_{\mathrm{opt}}$, which now carries the
  $O(\sigma/\sqrt{TB})$-type noise floor of SGD in addition to the deterministic,
  conditioning-driven contraction. In the convex setting one-pass SGD is known to attain
  the optimal statistical rates, so nothing is lost in the exchange.
\end{enumerate}

\section{The experiment: fixed budget, varying $n$}

Fix the model, the loss (hence the conditioning), and a compute budget of $T$ forward
passes. Train on datasets of size $n\in\{n_1<n_2<\dots\}$ plus one streaming run
($n=\infty$), and report the relative $L^2$ error on a fixed held-out test set. The
decomposition \eqref{eq:decomp} predicts the shape of the curve:
\begin{equation}
  \mathrm{err}(n)
  \;\approx\;
  \underbrace{\eps_{\mathrm{app}} + \eps_{\mathrm{opt}}(T)}_{\text{$n$-independent floor}}
  \;+\;
  \underbrace{O\bigl((C_\Theta/n)^{\alpha}\bigr)}_{\text{estimation}} ,
\end{equation}
i.e.\ a monotonically decreasing curve that \emph{saturates} at the value attained by
the streaming run. The crossover scale $n^\star$---where the estimation term drops below
the floor---marks the dataset size beyond which more data is useless \emph{at this
budget}; for $n\ge n^\star$ the finite-data runs should be statistically
indistinguishable from streaming. The figure then makes a clean claim: at fixed
optimization budget, the infinite-data limit is (weakly) optimal, and physics-informed
training gets it for free (Remark~\ref{rem:free}).

\begin{remark}[Budget accounting]
``Same number of forward passes'' is the right currency, because per-sample cost is
identical across runs (the residual assembly is negligible next to the FNO forward). For
small $n$ the optimizer will plateau---or overfit---long before exhausting $T$; early
stopping on a \emph{validation} set (patience of $x$ evaluations) then truncates the run.
This is fair, not generous, to small $n$: the stopped run reports the best risk it ever
achieved, which is exactly what a practitioner would extract from that dataset.
\end{remark}

\begin{remark}[Design caveats]
\leavevmode
\begin{enumerate}
  \item \textbf{Never stop on the test set.} Early stopping must monitor a validation
  set; the test set is touched once per reported number. Stopping on the test error
  biases the reported minimum downward, most severely for the small-$n$ runs whose
  validation curves are noisiest---precisely the runs the figure is about.
  \item \textbf{Seeds.} The small-$n$ runs inherit variance from the draw of $D_n$
  itself; repeat over a few dataset seeds and plot mean $\pm$ spread, otherwise the
  left end of the curve is a single noisy realization.
  \item \textbf{Test-set resolution.} A test set of size $m$ estimates the population
  test error to $O(m^{-1/2})$ relative accuracy; $m$ must be large enough that this is
  below the gaps the figure claims to resolve (a few hundred samples suffices at the
  effect sizes we expect).
  \item \textbf{Plotting $n=\infty$.} On a logarithmic $n$-axis the streaming run has no
  abscissa; draw it as a horizontal reference line (the floor), which is also the more
  honest visual statement.
\end{enumerate}
\end{remark}

\section*{Reference}

L.~Bottou and O.~Bousquet, \emph{The tradeoffs of large scale learning}, Advances in
Neural Information Processing Systems~20 (2008).

\end{document}
